\chapter{历年真题索引与模板映射}\label{ch:11}

\section{怎样查题目}\label{sec:11-1}
这一章可以按场次、题号查找，也可以按后面的题面关键词反查。前面列出近两年的题目，后面按 A、B、C、D/E 分类整理历年题目。

模板编号指向能借用的代码片段。例如 LCA 可以帮助找树上路径，却不会直接算出路径 mex；区间加线段树也不能原样用来做区间赋值。练习时仍要结合题目的条件。

\section{按场次查找}\label{sec:recent}
这里收录 2024 年 9 月至 2026 年 5 月的八场题目，共 40 题。题号链接到公开题面，表中列出可参考的写法和容易漏掉的条件。

各场起始页：202409（第\pageref{recent:202409}页）、202412（第\pageref{recent:202412}页）、202503（第\pageref{recent:202503}页）、202506（第\pageref{recent:202506}页）、202509（第\pageref{recent:202509}页）、202512（第\pageref{recent:202512}页）、202603（第\pageref{recent:202603}页）、202605（第\pageref{recent:202605}页）。

\Needspace{20\baselineskip}
\section{202409 场次速查}\label{recent:202409}
\begingroup
\small
\setlength{\tabcolsep}{4pt}
\begin{longtable}{p{3.1cm}p{5.5cm}p{6.2cm}}
\toprule
题号与题名 & 可参考的做法 & 要留意的条件 \\
\midrule
\endfirsthead
\toprule
题号与题名 & 可参考的做法 & 要留意的条件 \\
\midrule
\endhead
\href{https://oj.shumeng.tech/p/CSP202409A}{202409-1} 密码 & 逐字符分类与频次统计；\Tref{106}。 & 字母是一类，大小写字符分别计数；先检查三类字符，再判断频次。 \\
\href{https://oj.shumeng.tech/p/CSP202409B}{202409-2} 字符串变换 & 有限字符映射重复执行；\Tref{113} 的倍增或逐环取模。 & 含空格，用 getline；每个询问从原串出发；未给出的映射保持自身。 \\
\href{https://oj.shumeng.tech/p/CSP202409C}{202409-3} 补丁应用 & 按行解析块、校验旧内容、搜索偏移并替换；\Tref{601}、\Tref{602}。 & 保留行首空格；先检验合法性，再修改。简化子任务：合法、无偏移、小文本。 \\
\href{https://oj.shumeng.tech/p/CSP202409D}{202409-4} 通讯延迟 & 节点与覆盖它的基站建二部图，边权分别为延迟/0；\Tref{204}。 & 小数据可按基站枚举点对建边；正方形含边界；不可达输出 Nan。 \\
\href{https://oj.shumeng.tech/p/CSP202409E}{202409-5} 木板切割 & 序列分裂后统计不同色与颜色段；先做小规模 vector 模拟。 & 段编号固定，不随切割重排；空木板两项均为 0。小数据 n,m,k 不超过 2000。 \\
\bottomrule
\end{longtable}
\endgroup
\Needspace{20\baselineskip}
\section{202412 场次速查}\label{recent:202412}
\begingroup
\small
\setlength{\tabcolsep}{4pt}
\begin{longtable}{p{3.1cm}p{5.5cm}p{6.2cm}}
\toprule
题号与题名 & 可参考的做法 & 要留意的条件 \\
\midrule
\endfirsthead
\toprule
题号与题名 & 可参考的做法 & 要留意的条件 \\
\midrule
\endhead
\href{https://oj.shumeng.tech/p/CSP202412A}{202412-1} 移动 & 逐命令尝试新坐标；第\ref{ch:05}章方向与边界。 & 越界指令被忽略；题目 f/b 改 y，l/r 改 x，不要套错坐标轴。 \\
\href{https://oj.shumeng.tech/p/CSP202412B}{202412-2} 梦境巡查 & 将各段最低初始能量写成前缀消耗，维护前后缀最大值。 & 去掉第 i 处补给，只会抬高它之后的要求；小数据逐个禁用补给重算。 \\
\href{https://oj.shumeng.tech/p/CSP202412C}{202412-3} 缓存模拟 & 每个缓存组独立 LRU，附有效位、脏位；\Tref{605} 提供顺序维护。 & 淘汰脏行先写回，再载入；命中读写都更新次序。单路或无需淘汰子任务可先做。 \\
\href{https://oj.shumeng.tech/p/CSP202412D}{202412-4} 跳房子 & 状态是回退后的格子；BFS 枚举可落点；\Tref{201} 提供框架。 & 一次跳到 j 后必须回退到 j-a[j]；n 不超过 1000 时可显式枚举边。 \\
\href{https://oj.shumeng.tech/p/CSP202412E}{202412-5} 梦魔 & 两侧扩张与最低初始攻击力；小数据先逐起点重算，需专门推导。 & 只有左右最近的梦魔能打；收益为非负时可优先打当前可击杀对象。更新后不能沿用旧答案。 \\
\bottomrule
\end{longtable}
\endgroup
\Needspace{20\baselineskip}
\section{202503 场次速查}\label{recent:202503}
\begingroup
\small
\setlength{\tabcolsep}{4pt}
\begin{longtable}{p{3.1cm}p{5.5cm}p{6.2cm}}
\toprule
题号与题名 & 可参考的做法 & 要留意的条件 \\
\midrule
\endfirsthead
\toprule
题号与题名 & 可参考的做法 & 要留意的条件 \\
\midrule
\endhead
\href{https://oj.shumeng.tech/p/CSP202503A}{202503-1} 数值积分 & 枚举闭区间内偶数 x，累加题设公式；第\ref{ch:05}章公式题。 & 这是题设离散估计，不是求解析积分；起始偶数、端点与最后乘 2。 \\
\href{https://oj.shumeng.tech/p/CSP202503B}{202503-2} 机器人饲养指南 & 完全背包：一天吃 i 个视为重量 i、价值 A[i]；\Tref{302}、\Tref{303}。 & 须喂完 n 个，使用恰好状态；收益未必单调，不按单位收益贪心。 \\
\href{https://oj.shumeng.tech/p/CSP202503C}{202503-3} 模板展开 & 直接赋值存长度值，间接赋值存表达式依赖；按需递归求长度并取模。 & 间接引用每次求值要反映当前变量；自引用直接赋值先算旧值；只求长度，无需构造巨串。 \\
\href{https://oj.shumeng.tech/p/CSP202503D}{202503-4} 集体锻炼 & 区间 gcd 的贡献求和；\Tref{512}。小数据双循环滚动更新 gcd。 & 贡献还乘左右端点 l、r，及时取模；满分需压缩相同 gcd 状态，未收专用模板。 \\
\href{https://oj.shumeng.tech/p/CSP202503E}{202503-5} 收费标准评估 & 带负权的最大连通子树；小数据更新后重跑树 DP。 & 只操作 1/2 的小数据先做；dp[u]=a[u]+各孩子 max(0,dp[v])，非空连通块不能全部取 0。 \\
\bottomrule
\end{longtable}
\endgroup
\Needspace{20\baselineskip}
\section{202506 场次速查}\label{recent:202506}
\begingroup
\small
\setlength{\tabcolsep}{4pt}
\begin{longtable}{p{3.1cm}p{5.5cm}p{6.2cm}}
\toprule
题号与题名 & 可参考的做法 & 要留意的条件 \\
\midrule
\endfirsthead
\toprule
题号与题名 & 可参考的做法 & 要留意的条件 \\
\midrule
\endhead
\href{https://oj.shumeng.tech/p/CSP202506A}{202506-1} 正态分布 & 只求表格行列，用整数 t=100(n-mu)/sigma；行 t/10+1，列 t\%10+1。 & 题设 sigma 是 100 的因子，整数计算可避免浮点截断误差；无需计算概率函数。 \\
\href{https://oj.shumeng.tech/p/CSP202506B}{202506-2} 机器人复健指南 & 马步八方向 BFS，统计距离不超过 k 的格子；\Tref{201}。 & 移动是 (1,2)/(2,1) 的带符号组合，不是邻接八格；起点算入答案且只计一次。 \\
\href{https://oj.shumeng.tech/p/CSP202506C}{202506-3} 消息解码 & 定长位字段拆解、字符集转换、最近历史映射；\Tref{513} 仅提供进制基础。 & 先用旧历史解本条，再更新历史；间接推断的代号不能回填成直接收到。先做指定消息类型子任务。 \\
\href{https://oj.shumeng.tech/p/CSP202506D}{202506-4} 月票发行 & 固定短模式的自动机/状态 DP；跟踪先 ccf 后 cspark。 & 短串子任务也不宜枚举所有 26 的 n 次方字符串；普通 KMP 只匹配，不能直接计数所有构造。 \\
\href{https://oj.shumeng.tech/p/CSP202506E}{202506-5} 博物馆 & 删除一个点后的连通块统计；小图枚举施工点，BFS 各块的攻略点数。 & 每个施工点取某连通块中最多的攻略点数再求和；施工点不能参观；大图需割点/点双等未收内容。 \\
\bottomrule
\end{longtable}
\endgroup
\Needspace{20\baselineskip}
\section{202509 场次速查}\label{recent:202509}
\begingroup
\small
\setlength{\tabcolsep}{4pt}
\begin{longtable}{p{3.1cm}p{5.5cm}p{6.2cm}}
\toprule
题号与题名 & 可参考的做法 & 要留意的条件 \\
\midrule
\endfirsthead
\toprule
题号与题名 & 可参考的做法 & 要留意的条件 \\
\midrule
\endhead
\href{https://oj.shumeng.tech/p/CSP202509A}{202509-1} 蒙特卡洛 & 统计给定点满足 x*x+y*y<=a*a，输出 4.0*count/n。 & 输入已经给出样本，不要自己随机生成；含圆周；两位小数可放大 100 后用整数平方比较。 \\
\href{https://oj.shumeng.tech/p/CSP202509B}{202509-2} 水印检查 & 枚举 5x9 位置；每块得到阈值区间，差分合并；\Tref{103}。 & 黑点要求 k 大于其最大灰度，白点要求 k 不超过其最小灰度；区间为 [maxBlack+1,minWhite]。 \\
\href{https://oj.shumeng.tech/p/CSP202509C}{202509-3} HTTP 头信息 & 静态表、动态表、三类指令与 Huffman 解码；\Tref{602}、\Tref{601}。 & 区分普通 H 转义与编码串，去掉末尾填充位。只含表引用，或无 Huffman 的子任务可先做。 \\
\href{https://oj.shumeng.tech/p/CSP202509D}{202509-4} 造题计划（上） & 树上路径 mex；小数据沿路径收集值再扫描最小缺失数；\Tref{208}。 & mex 从 0 开始；题设权值为排列；LCA 只帮找路径，不直接给 mex。 \\
\href{https://oj.shumeng.tech/p/CSP202509E}{202509-5} 造题计划（下） & 有先后限制的出题/验题选择；小 n 可按天做状态 DP。 & 状态记录已出/已验数量及最小精力；同天上午出完可下午验，任何时刻已验不超过已出。 \\
\bottomrule
\end{longtable}
\endgroup
\Needspace{20\baselineskip}
\section{202512 场次速查}\label{recent:202512}
\begingroup
\small
\setlength{\tabcolsep}{4pt}
\begin{longtable}{p{3.1cm}p{5.5cm}p{6.2cm}}
\toprule
题号与题名 & 可参考的做法 & 要留意的条件 \\
\midrule
\endfirsthead
\toprule
题号与题名 & 可参考的做法 & 要留意的条件 \\
\midrule
\endhead
\href{https://oj.shumeng.tech/p/CSP202512A}{202512-1} 集合 & 比较真实集合相等与异或指纹相等是否一致；\Tref{106}、排序。 & 异或相同不能证明集合相同；输出的是该判断是否正确，不是集合是否相等。 \\
\href{https://oj.shumeng.tech/p/CSP202512B}{202512-2} 数字变换 & 输入值域仅 512，枚举全部输入算 F，再建立逆映射；\Tref{113} 的有限状态思想。 & 先按题图精确实现三段位变换；复杂度约 512*m+n，无须对每个查询重跑全部操作。 \\
\href{https://oj.shumeng.tech/p/CSP202512C}{202512-3} 图片解码 & 密钥倒序，逆转局部旋转/翻转和整体旋转；第\ref{ch:07}章矩阵模拟。 & 逆序时先撤销整体旋转，再撤销局部；Z 小或只有翻转的子任务先逐格做。去掉补齐问号。 \\
\href{https://oj.shumeng.tech/p/CSP202512D}{202512-4} C 形阵 & 由等式得到 F=B，枚举 C,E 为 B 平方的约数，再算 A,G,D；\Tref{512}。 & 小 n 枚举并验证整除、所有等式与不同元素数；完美要求恰好 6 种，不是 7 种。 \\
\href{https://oj.shumeng.tech/p/CSP202512E}{202512-5} 数据抢修 & 可重集合划分为两两异或至少 W 的组；最小数据回溯分组并剪枝。 & 重复权值也占独立元素；合并后 v 离线；不要套只含小写字母的 Trie。大数据专用解法未收录。 \\
\bottomrule
\end{longtable}
\endgroup
\Needspace{20\baselineskip}
\section{202603 场次速查}\label{recent:202603}
\begingroup
\small
\setlength{\tabcolsep}{4pt}
\begin{longtable}{p{3.1cm}p{5.5cm}p{6.2cm}}
\toprule
题号与题名 & 可参考的做法 & 要留意的条件 \\
\midrule
\endfirsthead
\toprule
题号与题名 & 可参考的做法 & 要留意的条件 \\
\midrule
\endhead
\href{https://oj.shumeng.tech/p/CSP202603A}{202603-1} 平衡数 & 逐位统计正整数二进制中的 0/1；第\ref{ch:05}章位运算。 & 不统计最高 1 之前的补零；重复输入分别计数。题名是平衡数。 \\
\href{https://oj.shumeng.tech/p/CSP202603B}{202603-2} 机器人项目管理 & 普通任务 01 背包，灵活任务按收益/咖啡排序；枚举预算分配；\Tref{301}。 & 只对可分割的灵活任务用单位收益贪心；普通任务不能拆；结果保留足够小数。 \\
\href{https://oj.shumeng.tech/p/CSP202603C}{202603-3} 进程通信 & 进程接口、队列与内存区间模拟；\Tref{601}、\Tref{602}、第\ref{ch:07}章状态设计。 & 分开维护占用和有对象两种状态；最优适应先最短再最左。小范围先逐地址/区间模拟。 \\
\href{https://oj.shumeng.tech/p/CSP202603D}{202603-4} 异或 & k 进制逐位不进位加法、f 值与区间操作；先做小范围直接模拟。 & k 是奇数，不能直接用 C++ 二进制异或符号代替；\Tref{402} 的普通加法懒标记也不能直接套。 \\
\href{https://oj.shumeng.tech/p/CSP202603E}{202603-5} 旅游计划 & 树上路径和道路修复；小图逐计划走唯一路径，模拟备胎与维修。 & 强制在线时先按规则解码输入；计划途中轮胎状态要重置；镜像分 Easy/Hard，按原题约束区分。 \\
\bottomrule
\end{longtable}
\endgroup
\Needspace{20\baselineskip}
\section{202605 场次速查}\label{recent:202605}
\begingroup
\small
\setlength{\tabcolsep}{4pt}
\begin{longtable}{p{3.1cm}p{5.5cm}p{6.2cm}}
\toprule
题号与题名 & 可参考的做法 & 要留意的条件 \\
\midrule
\endfirsthead
\toprule
题号与题名 & 可参考的做法 & 要留意的条件 \\
\midrule
\endhead
\href{https://oj.shumeng.tech/p/CSP202605A}{202605-1} 银行家舍入 & 按字符串拆整数部分 a 与一位小数 b；第\ref{ch:04}章舍入说明。 & b=5 时向偶数取整，其余同四舍五入；不要用 llround 当成银行家舍入。 \\
\href{https://oj.shumeng.tech/p/CSP202605B}{202605-2} 机器人宿管指南 & 二分机器人数量，逐天模拟先变质后吃苹果；\Tref{107}、\Tref{510}。 & 每天损耗按当日剩余数量向上取整；第 m 天吃完恰为 0 仍可行；中间乘法用 long long。 \\
\href{https://oj.shumeng.tech/p/CSP202605C}{202605-3} 死锁优化 & 进程状态机，分申请/收益/释放三阶段；\Tref{603}、\Tref{110}。 & 普通、A/B/C 类分别处理；同段申请与强占优先级按题面。先完成无特殊进程子任务并识别死锁。 \\
\href{https://oj.shumeng.tech/p/CSP202605D}{202605-4} 石子游戏 & 胜负公式已给出；小区间枚举最后一段做划分 DP。 & 奇数位是相对子段起点的奇数位；预处理奇偶前缀异或，dp[r]=max(dp[l-1]+win(l,r))。 \\
\href{https://oj.shumeng.tech/p/CSP202605E}{202605-5} 绝世好串 & 树上聚拢与字典序；极小 n 枚举合法聚拢和拼接顺序作验证。 & 第二次起必须参与已有长度至少 2 的串；随意排序所有字符不一定可实现。满分专用解法未收录。 \\
\bottomrule
\end{longtable}
\endgroup


\section{历年整场索引（201912--202406）}

下表按场次汇总题目。各题的具体入口见后面的 A、B、C 和 D/E 索引。

\begingroup
\small
\setlength{\tabcolsep}{3pt}
\begin{longtable}{p{1.55cm}p{2.8cm}p{2.8cm}p{3.35cm}p{4.6cm}}
\toprule
场次 & A/B 题 & C 题 & D/E 题信号 & 判断树重点 \\
\midrule
\endfirsthead
\toprule
场次 & A/B 题 & C 题 & D/E 题信号 & 判断树重点 \\
\midrule
\endhead
202406 & 矩阵重塑（一/二） & 文本分词 & 货物调度、哥德尔机 & 矩阵下标、文本 token、优先队列、贪心/背包。 \\
202403 & 词频统计、相似度计算 & 化学方程式配平 & 十滴水、文件夹合并 & 计数、集合交并、表达式/方程、set/堆、DFS 序/线段树。 \\
202312 & 仓库规划、因子化简 & 树上搜索 & 宝藏、彩色路径 & 暴力枚举、质因数分解、树 DFS、前缀/分块、状压 DP。 \\
202309 & 坐标变换（一/二） & 梯度求解 & 阴阳龙、阻击 & 坐标前缀合成、后缀表达式、set、树形 DP。 \\
202305 & 重复局面、矩阵运算 & 解压缩 & 电力网络、闪耀巡航 & 字符串哈希、矩阵乘法顺序、位运算解析、最短路/状压。 \\
202303 & 田地丈量、垦田计划 & LDAP & 星际网络II、施肥 & 矩形相交、二分答案、递归表达式、线段树/离散化。 \\
202212 & 现值计算、训练计划 & JPEG 解码 & 聚集方差、星际网络 & 公式、拓扑顺序、矩阵扫描、启发式合并、线段树建图。 \\
202209 & 如此编码、何以包邮 & 防疫大数据 & 通信/调度类复杂题 & 进制分解、01 背包、时间窗口、集合维护。 \\
202206 & 归一化处理、寻宝大冒险 & 角色授权 & 高级模拟/几何题 & 小数输出、坐标哈希、权限集合、规则匹配。 \\
202203 & 未初始化警告、出行计划 & 计算资源调度器 & 图/数据结构题 & set 判断、差分区间、对象筛选、优先级排序。 \\
202112 & 序列查询、新解 & 登机牌号码 & 磁盘文件操作、极差路径 & 分段贡献、数学优化、编码解析、线段树。 \\
202109 & 数组推导、非零段划分 & 脉冲神经网络 & 收集卡牌、箱根山岳险天下 & 前缀最大反推、贡献扫描、大模拟、概率/图。 \\
202104 & 灰度直方图、邻域均值 & DHCP服务器 & 校门外的树、疫苗运输 & 计数数组、二维前缀、大模拟状态机、DP。 \\
202012 & 安全指数、最佳阈值 & 带配额的文件系统 & 食材运输、星际旅行 & 公式、排序阈值扫描、目录树/配额、大模拟到树/图。 \\
202009 & 检测点查询、风险人群筛查 & 点亮数字人生 & 星际旅行、密信与计数 & 距离排序、连续区间判断、拓扑/逻辑门、图与计数。 \\
202006 & 线性分类器、稀疏向量 & Markdown渲染器 & 1246、乔乔和牛牛逛超市 & 坐标代入、双指针、文本渲染、DP/高级优化。 \\
201912 & 报数、回收站选址 & 化学方程式 & 区块链、魔数 & 简单模拟、坐标哈希、表达式解析、图/事件模拟。 \\
\bottomrule
\end{longtable}
\endgroup

练完一套题后，可以在相应场次旁记下自己漏看的条件或想错的做法。

\section{A 题索引}

\begingroup
\small
\setlength{\tabcolsep}{3pt}
\begin{longtable}{p{1.8cm}p{2.8cm}p{2.6cm}p{3.5cm}p{3.5cm}}
\toprule
编号 & 题名 & 题型 & 主要考点 & 翻到哪里 \\
\midrule
\endfirsthead
\toprule
编号 & 题名 & 题型 & 主要考点 & 翻到哪里 \\
\midrule
\endhead
202406-1 & 矩阵重塑（其一） & 下标 / 矩阵 & 一维二维下标互转 & A 题二维数组 \\
202403-1 & 词频统计 & 统计 & 循环计数、频率 & A 题频率统计 \\
202312-1 & 仓库规划 & 枚举 / 判断 & 双重循环、支配关系 & A 题简单模拟 \\
202309-1 & 坐标变换（其一） & 坐标模拟 & 坐标平移 & A 题坐标入门 \\
202305-1 & 重复局面 & 哈希统计 & 字符串 / 局面计数 & A 题频率统计、字符串 \\
202303-1 & 田地丈量 & 坐标 / 几何 & 矩形相交面积 & A 题坐标距离、公式 \\
202212-1 & 现值计算 & 公式 / 小数 & 幂、浮点计算 & A 题公式题 \\
202209-1 & 如此编码 & 进制 / 模拟 & 数位分解 & A 题进制处理 \\
202206-1 & 归一化处理 & 小数 / 统计 & 平均、方差、小数输出 & A 题小数输出 \\
202203-1 & 未初始化警告 & 模拟 / set & 变量出现判断 & A 题存在判断、set \\
202112-1 & 序列查询 & 模拟 / 查询 & 下标、区间含义 & A 题区间入门 \\
202109-1 & 数组推导 & 模拟 / 统计 & 最小最大和 & A 题简单统计 \\
202104-1 & 灰度直方图 & 统计 & 计数数组 & A 题计数数组 \\
202012-1 & 期末预测之安全指数 & 统计 / 公式 & 加权求和、取 max & A 题简单统计 \\
202009-1 & 称检测点查询 & 排序 / 距离 & 距离计算、排序 & A 题排序、坐标距离 \\
202006-1 & 线性分类器 & 模拟 / 几何判断 & 代入直线、分类统计 & A 题公式题、坐标 \\
201912-1 & 报数 & 模拟 & 跳过含 7 或 7 倍数 & A 题简单模拟 \\
201909-1 & 小明种苹果 & 统计 & 总量、最大减少量 & A 题简单统计 \\
201903-1 & 小中大 & 排序 / 输出格式 & 中位数、整数/小数输出 & A 题排序、格式 \\
201812-1 & 小明上学 & 时间模拟 & 红绿灯时间累计 & A 题时间处理 \\
201809-1 & 卖菜 & 数组模拟 & 邻域平均 & A 题数组边界 \\
201803-1 & 跳一跳 & 状态模拟 & 连续奖励计分 & A 题状态扫描 \\
201712-1 & 最小差值 & 排序 & 相邻差最小 & A 题排序 \\
201709-1 & 打酱油 & 数学 / 贪心 & 优惠规则分解 & A 题公式题 \\
201703-1 & 分蛋糕 & 贪心模拟 & 累加达到阈值分组 & A 题简单模拟 \\
201612-1 & 中间数 & 排序 / 统计 & 比它大和小数量 & A 题排序/计数 \\
201609-1 & 最大波动 & 扫描 & 相邻差绝对值最大 & A 题简单统计 \\
201604-1 & 折点计数 & 扫描 & 左右邻居比较 & A 题简单扫描 \\
201512-1 & 数位之和 & 数位 & 十进制拆位 & A 题数字处理 \\
201509-1 & 数列分段 & 扫描 & 相邻变化计数 & A 题简单扫描 \\
201503-1 & 图像旋转 & 矩阵 & 下标变换输出 & A 题二维表格 \\
201412-1 & 门禁系统 & 统计 & 出现次数累计 & A 题频率统计 \\
201409-1 & 相邻数对 & 排序 / 计数 & 排序后相邻比较 & A 题排序 \\
201403-1 & 相反数 & 统计 / 查找 & 取相反数、集合判断 & A 题存在判断 \\
201312-1 & 出现次数最多的数 & 统计 & 计数数组 / map & A 题频率统计 \\
\bottomrule
\end{longtable}
\endgroup

\section{B 题索引}

\begingroup
\small
\setlength{\tabcolsep}{3pt}
\begin{longtable}{p{1.8cm}p{2.8cm}p{2.6cm}p{3.5cm}p{3.5cm}}
\toprule
编号 & 题名 & 题型 & 主要考点 & 翻到哪里 \\
\midrule
\endfirsthead
\toprule
编号 & 题名 & 题型 & 主要考点 & 翻到哪里 \\
\midrule
\endhead
202406-2 & 矩阵重塑（其二） & 矩阵 / 模拟 & 下标映射、批量输出 & B 题矩阵、复杂度 \\
202403-2 & 相似度计算 & 集合 & 交集、并集、去重 & B 题哈希统计 \\
202312-2 & 因子化简 & 数学 / 因数分解 & 质因数分解、幂 & B 题基础数学 \\
202309-2 & 坐标变换（其二） & 前缀 / 坐标变换 & 操作前缀合成 & B 题前缀预处理 \\
202305-2 & 矩阵运算 & 矩阵 / 模拟优化 & 计算顺序、复杂度 & B 题表格矩阵、复杂度 \\
202303-2 & 垦田计划 & 二分答案 / 贪心 & 最小时间、check & B 题二分答案 \\
202212-2 & 训练计划 & 拓扑 / 模拟 & 最早最晚时间 & B 题拓扑排序入口 \\
202209-2 & 何以包邮？ & DP / 背包 & 01 背包 & B 题简单 DP \\
202206-2 & 寻宝！大冒险！ & 坐标 / 哈希 & 点集匹配 & B 题哈希统计、坐标 \\
202203-2 & 出行计划 & 差分 / 区间统计 & 时间区间覆盖 & B 题差分 \\
202112-2 & 序列查询新解 & 前缀 / 数学 & 分段贡献 & B 题前缀和、排序扫描 \\
202109-2 & 非零段划分 & 差分 / 扫描 & 贡献、扫描统计 & B 题差分、排序扫描 \\
202104-2 & 邻域均值 & 二维前缀和 & 子矩阵平均 & B 题二维前缀和 \\
202012-2 & 期末预测之最佳阈值 & 排序 / 前缀统计 & 阈值扫描、计数 & B 题排序加扫描 \\
202009-2 & 风险人群筛查 & 模拟 / 坐标 & 连续停留判断 & B 题排序扫描、模拟 \\
202006-2 & 稀疏向量 & 双指针 / map & 有序向量点积 & B 题双指针、哈希统计 \\
201912-2 & 回收站选址 & 坐标 / 哈希 & 邻居存在、评分统计 & B 题哈希坐标 \\
201909-2 & 小明种苹果（续） & 统计 / 标记 & 掉落判断、连续组 & B 题统计扫描 \\
201903-2 & 二十四点 & 表达式 / 栈 & 运算优先级、求值 & C 题表达式入门 \\
201812-2 & 小明放学 & 时间 / 模拟 & 周期取模、信号灯 & B 题时间处理 \\
201809-2 & 买菜 & 区间重叠 & 时间段交集累计 & B 题区间扫描 \\
201803-2 & 碰撞的小球 & 模拟 & 位置更新、碰撞反向 & B 题模拟 \\
201712-2 & 游戏 & 队列模拟 & 报数淘汰 & B 题队列 \\
201709-2 & 公共钥匙盒 & 事件排序 & 还钥匙优先、时间排序 & C 题事件模拟 \\
201703-2 & 学生排队 & 序列模拟 & 移动元素 & B 题 vector 操作 \\
201612-2 & 工资计算 & 反推 / 枚举 & 税率分段、枚举答案 & B 题分段函数 \\
201609-2 & 火车购票 & 模拟 / 贪心 & 连续座位优先 & B 题模拟 \\
201604-2 & 俄罗斯方块 & 矩阵模拟 & 碰撞检测、落点 & B 题矩阵模拟 \\
201512-2 & 消除类游戏 & 二维扫描 & 标记后统一清除 & B 题矩阵模拟 \\
201509-2 & 日期计算 & 日期 & 月份天数、闰年 & A/B 日期模板 \\
201503-2 & 数字排序 & 排序 / 频率 & 频次降序、值升序 & B 题排序比较器 \\
201412-2 & Z字形扫描 & 矩阵遍历 & 对角线方向切换 & B 题矩阵扫描 \\
201409-2 & 画图 & 二维标记 & 矩形覆盖面积 & B 题二维表格 \\
201403-2 & 窗口 & 模拟 / 栈序 & 点击命中、置顶 & B 题模拟、列表 \\
201312-2 & ISBN号码 & 字符串 / 校验 & 加权求和、校验码 & A 题字符串、公式 \\
\bottomrule
\end{longtable}
\endgroup

\section{C 题索引：大模拟与中档算法入口}

C 题按输入中的对象、规则和字符串格式查找。表中的做法可以作参考，具体实现还要看数据范围。

\begingroup
\small
\setlength{\tabcolsep}{3pt}
\begin{longtable}{p{1.75cm}p{2.55cm}p{3.0cm}p{3.6cm}p{3.0cm}}
\toprule
编号 & 题名 & 第一识别信号 & 模板入口 & 易错提醒 \\
\midrule
\endfirsthead
\toprule
编号 & 题名 & 第一识别信号 & 模板入口 & 易错提醒 \\
\midrule
\endhead
202406-3 & 文本分词 & 词表、token、合并规则 & \Tref{601}、\Tref{604}、字符串解析 & 输入文本和词表顺序不要混。 \\
202403-3 & 化学方程式配平 & 方程、元素、系数未知 & 表达式解析、线性方程 & 元素名大小写和括号倍数。 \\
202312-3 & 树上搜索 & 树、子树、搜索顺序 & DFS、树模拟 & 子树权值和、删除/保留状态。 \\
202309-3 & 梯度求解 & 表达式、变量、后缀计算 & 后缀栈求导 & 操作数顺序、取模和导数规则。 \\
202305-3 & 解压缩 & 编码串、块、长度 & 字符串解析、位运算 & 十六进制、长度字段、边界。 \\
202303-3 & LDAP & 递归条件表达式 & 递归下降、集合/bitset & 与/或/非优先级和括号。 \\
202212-3 & JPEG 解码 & 8x8 矩阵、Z 字形、变换 & 矩阵模拟、浮点/取整 & 下标顺序、四舍五入、截断范围。 \\
202209-3 & 防疫大数据 & 时间窗口、风险地区、人员轨迹 & \Tref{601}、set/map & 生效时间和过期时间。 \\
202206-3 & 角色授权 & 用户、角色、权限、资源 & map/set、规则匹配 & 通配符、用户/用户组关联、权限匹配。 \\
202203-3 & 计算资源调度器 & 节点选择、多级约束 & \Tref{601}、排序、set & 优先级和约束过滤顺序。 \\
202112-3 & 登机牌号码 & 编码、校验、字符串转换 & 进制/编码、模拟 & 字符集映射和补位。 \\
202109-3 & 脉冲神经网络 & 多轮状态更新 & 大模拟、滚动数组 & 同一轮更新要先累计后生效。 \\
202104-3 & DHCP服务器 & 请求/响应、租约状态 & \Tref{603} 状态机、事件时间 & 到期时间和状态转移顺序。 \\
202012-3 & 带配额的文件系统 & 目录树、文件大小、配额 & \Tref{606}、树、map & 路径创建失败要回滚。 \\
202009-3 & 点亮数字人生 & 逻辑门、依赖顺序 & \Tref{205} 拓扑排序 & 环检测、多个输入输出。 \\
202006-3 & Markdown渲染器 & 文本块、列表、段落 & \Tref{601}、字符串解析 & 空行分块和宽度换行。 \\
201912-3 & 化学方程式 & 化学式、元素守恒 & 字符串解析、map & 多位系数、括号嵌套。 \\
201909-3 & 字符画 & 像素块、RGB、转义输出 & 矩阵模拟、字符串输出 & 转义序列和连续颜色压缩。 \\
201903-3 & 损坏的RAID5 & 磁盘、条带、异或恢复 & 模拟、位运算 & 十六进制转换和缺盘数量。 \\
201812-3 & CIDR合并 & IP/前缀、排序、合并 & \Tref{513}、排序扫描 & IP 转整数、包含关系。 \\
201809-3 & 元素选择器 & CSS 选择器、层级匹配 & 字符串解析、树/栈 & 大小写、祖先关系。 \\
201803-3 & URL映射 & 路由规则、参数提取 & 字符串匹配、类型参数检查 & 参数类型、匹配失败。 \\
201712-3 & Crontab & 时间范围、周期触发 & \Tref{514}、事件枚举 & 星期、月份、闭区间。 \\
201709-3 & JSON查询 & 键值、嵌套对象、查询 & 字符串解析、栈/map & 转义字符和层级路径。 \\
201703-3 & Markdown & 标题、列表、强调 & 字符串解析 & 块级规则优先于行内规则。 \\
201612-3 & 权限查询 & 用户/角色/权限等级 & map/set、规则判断 & 通配等级和最大权限。 \\
201609-3 & 炉石传说 & 角色、随从、攻击 & \Tref{603} 状态机 & 死亡结算顺序。 \\
201604-3 & 路径解析 & 当前目录、相对路径 & \Tref{606} 路径规范化 & 多个斜杠、\Code{..} 到根目录。 \\
201509-3 & 模板生成系统 & 变量替换、文本模板 & \Tref{501} split、map & 引号内容、未定义变量为空。 \\
201403-3 & 命令行选项 & 选项、参数、命令行 & \Tref{602} 命令分发 & 有参/无参选项区别。 \\
\bottomrule
\end{longtable}
\endgroup

\section{D/E 题索引与小数据做法}

D/E 题先看有无适合暴力的小数据，或链、树、无修改等特殊条件。表中列出一些可尝试的方向。

\begingroup
\small
\setlength{\tabcolsep}{3pt}
\begin{longtable}{p{2.0cm}p{3.0cm}p{4.0cm}p{4.9cm}}
\toprule
编号 & 题名 & 参考方向 & 小数据做法 \\
\midrule
\endfirsthead
\toprule
编号 & 题名 & 参考方向 & 小数据做法 \\
\midrule
\endhead
202403-4 & 十滴水 & set/map、优先队列、连锁反应 & 小数据暴力模拟每次爆裂。 \\
202403-5 & 文件夹合并 & DFS 序、链表、线段树 & 只做小树 DFS 合并，或无修改子任务。 \\
202312-4 & 宝藏 & 分块、前缀统计 & 暴力枚举查询，小范围预处理。 \\
202312-5 & 彩色路径 & 状压 DP、折半搜索 & 颜色数小就状压，否则 DFS 小图。 \\
202309-4 & 阴阳龙 & 动态 set / 平衡树 & 小数据每次排序或线性扫描。 \\
202309-5 & 阻击 & 树形 DP、线段树、树剖 & 链/小树暴力 DFS。 \\
202305-4 & 电力网络 & 图论、最短路/割点信号 & 小图建图后暴力枚举。 \\
202305-5 & 闪耀巡航 & 最短路、状压 DP & 点数小用 Floyd + 状压。 \\
202303-4 & 星际网络II & 线段树、离散化 & 小范围直接数组模拟。 \\
202303-5 & 施肥 & 分治、线段树/树状数组 & 差分可拿部分分。 \\
202212-4 & 聚集方差 & 启发式合并 & 小数据暴力维护集合。 \\
202212-5 & 星际网络 & 线段树建图、高精度 & 小图建边后跑最短路。 \\
202112-4 & 磁盘文件操作 & 线段树、区间覆盖 & 小范围数组模拟。 \\
202104-4 & 校门外的树 & DP & 先写 $O(n^2)$ 或按区间暴力。 \\
202009-4 & 星际旅行 & 图/几何/矩阵类优化 & 小数据按定义模拟。 \\
201812-4 & 数据中心 & 最小生成树 & \Tref{207} Kruskal。 \\
201412-4 & 最优灌溉 & 最小生成树 & \Tref{207} Kruskal。 \\
\bottomrule
\end{longtable}
\endgroup

\section{题目反查模板索引}

如果你记得“题看起来像什么”，但想不起算法名，按这一表反查。

\begin{longtable}{p{4.2cm}p{5.0cm}p{5.0cm}}
\toprule
题面样子 & 优先翻看 & 典型真题 \\
\midrule
\endfirsthead
\toprule
题面样子 & 优先翻看 & 典型真题 \\
\midrule
\endhead
一堆对象，每个对象有多个属性，要求按规则选择 & C 题状态机、大模拟骨架、排序比较器 & 计算资源调度器、DHCP服务器、文本分词 \\
路径、目录、文件、配额、URL、JSON、Markdown & 字符串解析、路径规范化、递归结构 & 路径解析、带配额的文件系统、JSON查询、URL映射 \\
矩阵、图像、8x8、Z 字形、旋转、重塑 & 二维数组、矩阵扫描、下标变换 & JPEG 解码、矩阵重塑、Z字形扫描、图像旋转 \\
多次询问，每次问区间/矩形/时间覆盖 & 前缀和、差分、线段树 & 邻域均值、出行计划、磁盘文件操作 \\
要求最小时间、最大收益、阈值最优，且答案越大越容易/越难 & 二分答案 & 垦田计划、资源压缩题 \\
依赖关系、先后顺序、逻辑门、任务计划 & 拓扑排序 & 训练计划、点亮数字人生 \\
集合合并、连通块、最小连接代价 & 并查集、Kruskal & 数据中心、最优灌溉 \\
字符串反复变换、替换函数、循环节 & 映射图、环、倍增/取模 & 字符串变换类题 \\
\bottomrule
\end{longtable}

\section{按题面信号反查真题}

记不清题号时，可以按题面中的操作、对象或数据结构查找。

\begingroup
\small
\setlength{\tabcolsep}{3pt}
\begin{longtable}{p{3.0cm}p{3.6cm}p{4.2cm}p{4.0cm}}
\toprule
题面信号 & 第一判断 & 类似真题 & 模板入口 \\
\midrule
\endfirsthead
\toprule
题面信号 & 第一判断 & 类似真题 & 模板入口 \\
\midrule
\endhead
只要求统计数量、最大最小、频次 & A 题扫描/计数 & 出现次数最多的数、灰度直方图、词频统计 & A 题统计、\Tref{106} \\
相邻元素、连续段、分段数量 & 一遍扫描或排序后扫描 & 数列分段、折点计数、非零段划分 & A 题扫描、B 题排序扫描 \\
矩形、坐标点、覆盖面积 & 坐标公式或哈希点集 & 田地丈量、画图、回收站选址、寻宝大冒险 & 坐标、\Tref{106}、\Tref{109} \\
矩阵重塑、旋转、Z 字形、局部平均 & 二维下标或二维前缀 & 图像旋转、Z字形扫描、邻域均值、矩阵重塑 & \Tref{102}、矩阵扫描 \\
时间段、红绿灯、租约、到期 & 时间模拟或事件排序 & 小明上学、小明放学、公共钥匙盒、DHCP服务器 & \Tref{110}、\Tref{514}、\Tref{603} \\
稀疏数据，两列编号和值 & map 或双指针 & 稀疏向量、坐标变换、回收站选址 & \Tref{106}、\Tref{108} \\
阈值、最优分界点、排序后评价 & 排序扫描或前缀统计 & 最佳阈值、非零段划分 & \Tref{105}、\Tref{101} \\
最小时间、最大可行值、资源压缩 & 二分答案 & 垦田计划、资源压缩类题 & \Tref{107} \\
选择若干物品，金额达到目标 & 可达金额 DP & 何以包邮 & \Tref{303} \\
依赖任务、前置关系、逻辑门 & 拓扑排序 & 训练计划、点亮数字人生 & \Tref{205} \\
路径、目录、文件、配额 & 字符串解析 + 树形结构 & 路径解析、带配额的文件系统、文件夹合并 & \Tref{606}、\Tref{601}、\Tref{402} \\
JSON、URL、Markdown、模板替换 & 格式解析 & JSON查询、URL映射、Markdown、模板生成系统 & \Tref{501}、\Tref{607}、\Tref{601} \\
化学式、表达式、括号、方程 & 按语法和任务区分 & 化学方程式、配平、梯度求解、二十四点 & 算术见 \Tref{610}；化学/求导需专用实现 \\
缓存、最近最少使用、淘汰 & 链表/map 或队列模拟 & 缓存类模拟题 & \Tref{605} \\
对象有状态，命令改变状态 & 状态机 & DHCP服务器、炉石传说、角色授权、计算资源调度器 & \Tref{603}、\Tref{602} \\
每次取最优对象，可能删除失效对象 & 堆或 set，必要时懒删除 & 文本分词、任务调度类题、十滴水 & \Tref{604} \\
树上子树、祖先、路径、合并 & DFS 序、LCA、树形 DP & 树上搜索、文件夹合并、阻击 & \Tref{208}、\Tref{402} \\
动态区间修改查询 & 线段树/树状数组 & 磁盘文件操作、星际网络II、施肥 & \Tref{401}、\Tref{402} \\
最小连接代价、网络建设 & 最小生成树 & 最优灌溉、数据中心 & \Tref{207} \\
点少、状态多、颜色/集合限制 & 状压 DP & 彩色路径、闪耀巡航 & 第 8 章状压入口 \\
\bottomrule
\end{longtable}
\endgroup

\section{练习后怎样复盘}

做完一道题，可以对照最初的思路和实际做法，记下差别。

\begin{longtable}{p{3.0cm}p{5.2cm}p{5.4cm}}
\toprule
检查项 & 如果答不上来 & 应补到哪里 \\
\midrule
\endfirsthead
\toprule
检查项 & 如果答不上来 & 应补到哪里 \\
\midrule
\endhead
这题第一眼信号是什么？ & 只记得题名，不知道题面特征 & 本章“按题面信号反查真题”。 \\
数据范围卡掉了哪些暴力？ & 没估复杂度就写代码 & 第 3 章数据范围判断树。 \\
满分模板是什么？ & 只会题解，不会迁移 & 第 9 章模板库和第 3 章关键词索引。 \\
有没有更简单部分分？ & D/E 空着没写 & 第 8 章 D/E 部分分策略。 \\
我误判成了什么？ & 例如把二分题当贪心，把表达式当普通字符串 & 第 3 章容易误判分支。 \\
这题有没有固定易错点？ & 样例过了但提交错 & 第 10 章易错点总表。 \\
\bottomrule
\end{longtable}

\section{A/B 题型归纳}

\subsection{A 题高频题型}

\begin{longtable}{p{4cm}p{5cm}p{5cm}}
\toprule
题型 & 常见信号 & 对应章节 \\
\midrule
\endfirsthead
\toprule
题型 & 常见信号 & 对应章节 \\
\midrule
\endhead
计数统计 & 出现次数、直方图、满足条件数量 & A 题简单统计、频率统计 \\
公式计算 & 安全指数、现值、加权求和 & A 题公式题 \\
坐标处理 & 检测点、坐标变换、矩形 & A 题坐标和距离 \\
字符串/哈希 & 重复局面、字符串计数 & A 题字符串、频率统计 \\
二维数组 & 图像旋转、矩阵重塑、卖菜 & A 题数组边界、二维表格 \\
简单模拟 & 按题意判断、逐项检查 & A 题简单模拟 \\
\bottomrule
\end{longtable}

\subsection{B 题高频题型}

\begin{longtable}{p{4cm}p{5cm}p{5cm}}
\toprule
题型 & 常见信号 & 对应章节 \\
\midrule
\endfirsthead
\toprule
题型 & 常见信号 & 对应章节 \\
\midrule
\endhead
前缀和 / 二维前缀和 & 多次区间或矩形查询 & B 题前缀和 \\
差分 & 时间区间覆盖、区间修改 & B 题差分 \\
排序扫描 & 阈值、排名、分段统计 & B 题排序加扫描 \\
哈希 / 集合 & 坐标点集、稀疏数据、相似度 & B 题哈希统计 \\
二分答案 & 最小化最大时间、资源压缩 & B 题二分答案 \\
背包 DP & 选择若干使总和达到要求 & B 题简单 DP \\
拓扑 / 依赖 & 任务计划、前置关系 & B 题拓扑排序 \\
矩阵模拟 & 重塑、扫描、局部统计 & B 题矩阵处理 \\
数学分解 & 因子、约数、质因数 & B 题基础数学 \\
\bottomrule
\end{longtable}
